\chapter{Executive Summary}

The \texttt{MU+MBAD7112+DEV\_01} course (Project Management) is
\textbf{not ready for publication}. It is a template shell in early
population state. Most content is placeholder text, all \texttt{about/}
metadata files are empty, and unpublished draft components duplicate
and contradict live components.

The two program PDFs provide the roadmap
(\textit{Online-Course-Material-Development.pdf}) and the QA baseline
(\textit{DIGITAL CONTENT QUALITY ASSURANCE RUBRIC TEMPLATE.pdf}). This
recommendation aligns to:

\begin{itemize}
  \item Roadmap Phases 0--9
  \item The RACI matrix and tool responsibilities in Section 2.3
  \item Appendix B (Branding Elements)
  \item The 9-category QA rubric, expanded to 14 categories with
        scoring bands and gate mapping
  \item Open edX Teak capabilities, including \texttt{course\_wide\_css}
\end{itemize}

\textbf{Key decisions:}
\begin{enumerate}
  \item A common CSS file, scoped to course content only, is uploaded
        via Studio $\to$ Files \& Uploads and applied via
        \texttt{course\_wide\_css} (Teak-supported).
  \item Authors work in Open edX Studio as primary authoring tool.
        GitHub holds exported course snapshots, issue tracking, and
        GitHub Pages previews.
  \item The QA rubric is expanded to 14 categories with a 0--5 scoring
        band, a minimum pass threshold of 4.0 weighted average, and
        no category below 3.
  \item All \texttt{about/} metadata must be populated before catalog
        visibility can be approved.
  \item The \texttt{video\_iframe} XBlock is not standard in Teak and
        must be migrated to standard \texttt{video} XBlocks or
        confirmed as a custom install.
  \item Multiple-choice items with all \texttt{correct="true"} choices
        will only award a point for the first correct choice in Teak.
        The pre-course self-assessment must be redesigned as a survey
        or poll XBlock.
\end{enumerate}